无线电干扰源的快速清除关系到通信安全，危险环境中的排查需要具备自主作业能力。题目设定机器狗在圆域内搜索未知干扰源，依靠有界误差测向与近距离光学探测完成任务。本文围绕区域确定、主动选点和多源协同清除建立模型并设计求解算法。

\textbf{针对问题一}，本文建立\textbf{半平面交与最小包围圆模型}，将各次示向度转化为前向角域，联立半平面约束确定定位区域。通过线性规划判定区域状态，枚举边界交点提取顶点，再计算最远顶点对与支撑圆，求得区域直径和最小覆盖半径。证明直径圆能够覆盖区域的充要条件为最小包围圆半径等于区域直径的一半，并构造可由三次测向得到的等边三角形反例，其直径为\(38\,\mathrm{m}\)，最小包围圆半径为\(21.94\,\mathrm{m}\)，说明直径小于\(40\,\mathrm{m}\)仍不足以保证清除，保证清除应以最小包围圆半径不超过\(20\,\mathrm{m}\)为判据。

\textbf{针对问题二}，本文建立\textbf{基于贝叶斯更新的加权选点模型}，联合首次角域、目标圆域及固定接收半径的约束，以主观概率描述未知量，预测第二次各类反馈及其剩余区域。将期望剩余半径与无需补测的完成概率归一化加权，采用分段高斯求积、网格搜索与多初值Nelder--Mead算法求解，汇集不同权重下的优选点形成候选区域。在首次点为原点、示向正东的算例中，提高完成概率权重时，优选点向首次点靠近，完成概率提高而期望剩余半径增大；两项指标等权时，优选点为\((813.8,\pm521.0)\,\mathrm{m}\)，期望剩余半径为\(25.67\,\mathrm{m}\)，完成概率为\(37.34\%\)。

\textbf{针对问题三}，本文建立\textbf{基于可行区域的搜索定位清除协同模型}，维护各频道的位置约束及覆盖记录，将扫描与目标处理共同安排到移动路线中，并按局部时间估计选择补测或提前试清。采用多起点近邻构造、路径片段反转与逐步重规划求解，以已知源全清及未知频道整网格排除作为结束依据。三次正式测试分别清除15、12、10个源，平均定位清除时间分别为\(195.399\)、\(276.748\)、\(248.719\,\mathrm{s}\)。两批官方演练共80场均全部清除，场均每源耗时分别为\(231.71\,\mathrm{s}\)和\(226.34\,\mathrm{s}\)；其中50场批次较独立演练的文献适配基线均值降低\(30.54\%\)，主要收益来自减少移动。

\textbf{针对问题四}，本文建立\textbf{朝向鲁棒搜索与联合可行域模型}，将源位置、接收半径、类型和朝向共同纳入观测约束，利用有信号与无信号历史收缩可行域。通过站点凸包与距离条件设计22站搜索布局，经4988个区域单元校核覆盖全部朝向；以有限假设评价测点，结合多初始路线与2-opt改进协调搜索和清除。所选策略在30场官方演练中全部清除390个源，场均每源耗时为\(458.37\,\mathrm{s}\)，未出现光学清除失败。
